% =====================================================================
%  Выводы по лабораторной работе.
% =====================================================================

\labpartbreak{Выводы}

В ходе выполнения лабораторной работы было разработано и практически
проверено веб-приложение, реализующее полный набор операций из задания
(регистрация, авторизация, CRUD и поиск для конфиденциальных и
неконфиденциальных данных, деавторизация), а также выполнен анализ его
защищённости по всем четырём требуемым направлениям: политике
безопасности операционной системы, политике безопасности компонентов
приложения и их взаимодействия, сетевой политике безопасности и
защищённости конфиденциальной информации. В отличие от первой
лабораторной работы, где контроль доступа обеспечивался средствами ОС и
СУБД, здесь все рубежи защиты -- от хеширования пароля до шифрования
конфиденциального поля -- реализованы и обоснованы на уровне кода самого
приложения, что и являлось целью работы.

Семантически работа показала, что защищённость приложения складывается из
нескольких независимых, не подменяющих друг друга механизмов: то, что
пароль нельзя восстановить из базы (хеширование), не отменяет
необходимости защищать сам токен сессии от кражи (\texttt{HttpOnly},
\texttt{Secure}, \texttt{SameSite}), а то, что токен защищён, не отменяет
необходимости проверять владельца записи на каждом запросе к данным --
все три меры проверялись по отдельности и ни одна не компенсирует
отсутствие другой. Отдельно подтверждено практикой, что «конфиденциальные
данные» -- это не декларация в описании поля, а измеримое свойство:
утверждение о шифровании было проверено не чтением кода, а прямым поиском
секретного значения в бинарном файле базы данных, и только этот
эксперимент, а не сам факт вызова функции \texttt{Encrypt}, можно считать
доказательством защищённости.

С прагматической точки зрения работа показала, что большая часть типовых
угроз из теоретической части задания (перебор, разглашение данных при
компрометации канала или носителя, использование уязвимостей в
коде) закрывается не экзотическими средствами, а корректным применением
стандартных, хорошо документированных механизмов платформы -- PBKDF2,
AES-GCM, JWT с флагами cookie, именованный CORS, параметризованные
запросы Entity Framework Core -- при условии, что каждый из них применён
осознанно и с понятным обоснованием, а не скопирован по шаблону. Также
показано, что автоматизированная проверка защищённости (скрипт
\texttt{verify.sh}, повторяющий сценарий атаки или обхода для каждого
рубежа защиты и утверждающий код ответа и содержимое файла базы) даёт
существенно более убедительное обоснование защищённости, чем ссылка на
код или документацию, и может быть воспроизведена при последующих
изменениях приложения для предотвращения регрессий безопасности.
